\chapter{Допуск общего каллиасова носителя}
\label{ch:v10-particle-wrinkle-dislocation-callias-carrier}

Профиль Каллиаса требует разделить пространственный спинор и внутреннюю
двойку массового скручивания. Минимальный условный носитель равен
\begin{equation}
 \mathcal H_{\rm C}
 =\mathbb C^2_{\rm spin}\otimes
  \mathbb C^2_{\rm twist}\otimes H_{15},
 \qquad \dim_{\mathbb C}\mathcal H_{\rm C}=60.
\end{equation}
Пространственные генераторы $\sigma_i\otimes I$ и twist-генераторы
$I\otimes\tau_a$ коммутируют. Обе тройки удовлетворяют точным соотношениям
Клиффорда, а проекторы знака массового профиля имеют ранги $30+30$.
Коэффициентный множитель $H_{15}$ сохраняет кратность индекса пятнадцать.

Накопленный носитель содержит необходимые ингредиенты по отдельности, но
не их требуемое произведение. Пространственный спинор вместе с $H_{15}$
имеет размерность $30$. KO6 particle/conjugate-пара не является внутренней
равнозарядной двойкой: недиагональный Pauli-переворот имеет коммутатор с
оператором заряда полного ранга $30$.

Двухвершинное клеточное ребро K43 даёт более интересную зацепку. На его
абстрактной $\mathbb C^2$ можно построить полный Pauli-триплет, коммутирующий
со скалярным зарядом. Однако унаследованное отображение этой двойки в
$\mathbb C^2_{\rm twist}\otimes H_{15}$ равно нулю и имеет ранг $0$.
Число $43$ также само по себе не задаёт равномерное усиление всех
пятнадцати физических каналов.

\begin{theorem}[Минимальная каллиасова архитектура согласована]
Носитель размерности $60$ с независимыми spin- и twist-факторами строго
реализует требуемые клиффордовы соотношения, два массовых сектора и
индексную кратность пятнадцать.
\end{theorem}
\begin{nogocriterion}[Real-удвоение не заменяет twist-двойку]
Сопряжённые particle/antiparticle-секторы имеют разные заряды, поэтому
Pauli-смешивание нарушает калибровочный коммутант. Клеточная двойка также
не становится twist-фактором без типизированного усиления над $H_{15}$.
\end{nogocriterion}
\begin{protocol}\begin{enumerate}
\item условная архитектура: \textbf{$10/10$};
\item отдельные унаследованные ингредиенты: \textbf{$5/5$};
\item происхождение равнозарядной twist-двойки: \textbf{$0/3$};
\item строгий допуск общего носителя: \textbf{$0/1$};
\item следующий гейт: \textbf{аудит источников twist-двойки}.
\end{enumerate}\end{protocol}
Машинный сертификат сохранён в
\path{s2t_v10_particle_wrinkle_dislocation_callias_profile_common_carrier_admission_gate_results.json}.